% GENERATED by tools/build_m05_code_sheet.py -- do not edit.
\begin{codebox}{Box 4}{The stochastic block model}
deg_corr = False
gt.seed_rng(1)
state = gt.minimize_blockmodel_dl(
    g, state_args=dict(deg_corr=deg_corr)
)
sbm = state.get_blocks().a
print(sbm)
pos = gt.graph_draw(
    g, pos=pos, vertex_text=g.vertex_index,
    vertex_fill_color=paint(sbm),
)
\end{codebox}
\begin{linenotes}
\item[1] This stores your choice of model under the name \texttt{deg\_corr}. \texttt{False} is the plain model of the lecture: the group alone decides how likely a friendship is, so the members of a group are alike in how many friends they have. \texttt{True} is the \emph{degree-corrected} model: every member may also have more or fewer friends than the others in the group. Change this one word later, to try the other model.
\item[2] The fit makes random choices. \texttt{seed\_rng(1)} fixes them, so that everyone in the room gets the same answer. Change the \texttt{1} and you may get a different grouping.
\item[3--5] \texttt{minimize\_blockmodel\_dl} fits the stochastic block model from the lecture. \texttt{dl} stands for \emph{description length}: of all the ways to group the nodes, it picks the one that describes the network in the fewest bits. It decides the number of groups too. \texttt{state\_args=dict(deg\_corr=deg\_corr)} hands your choice from line 1 to the model. The fitted model is stored as \texttt{state}.
\item[6] \texttt{state.get\_blocks()} gives the group of each node, in graph-tool's own array type. \texttt{.a} turns it into a plain array of numbers, like the list in Box 2.
\item[7] Print it: one number for each of the 34 nodes.
\item[8--11] The drawing call of Box 3, with the colours of the block model's groups. \texttt{pos=pos} reuses the positions from Box 3, so that each node stays in the same place and the pictures can be compared.
\end{linenotes}
